\documentclass{article}
\usepackage{graphicx}
\usepackage{amsmath}
\usepackage{amsfonts}
\usepackage{bookmark}
\usepackage{booktabs}
\usepackage{exercises}
\graphicspath{ {./images/5.3} }
\usepackage{indentfirst}
\usepackage{amsthm}
\usepackage{yhmath}
\usepackage{geometry}
\usepackage{placeins}
\usepackage{csquotes}
\usepackage{pgfplots}
\newcommand*{\oldneg}{\mathord{\sim}}
\newcommand\showdiv[1]{\overline{\smash{)}#1}}
\newcommand\Mydiv[1]{%
\smash{\raise.3ex\hbox{$\big)$}}\mkern-8mu
        \overline{\enspace\strut#1}}
\setlength{\tabcolsep}{5pt}
\theoremstyle{definition}
\newtheorem{definition}{Definition}

\setlength{\parindent}{.5in}
\setlength{\parskip}{12pt}
\renewcommand{\arraystretch}{2}
\geometry{margin=1in}
\newcommand{\?}{\stackrel{?}{=}}
\pgfplotsset{width=7cm,compat=1.18}\usepgfplotslibrary{statistics}

\title{5.5 Answers}
\begin{document}
\maketitle

\section*{1}
A fraction that contains a fraction in its numerator or denonminator is called a complex dfraction.

\section*{2}
Consider two rational expressions $\dfrac{a}{b}$ and $\dfrac{c}{d}$ with unlike denominators. They can be added by finding the least common denominator 
of the two fractions, which is the least common multiple of $b$ and $d$.

\section*{7}
\[ \dfrac{5x}{x+3}+\dfrac{15}{x+3}=\dfrac{5x+15}{x+3}=\dfrac{5(x+3)}{x+3}=5 \]

\section*{8}
\[ \dfrac{4x^2}{2x-1}-\dfrac{1}{2x-1}=\dfrac{4x^2-1}{2x-1}=\dfrac{(2x+1)(2x-1)}{2x-1}=2x+1 \]

\section*{13}
We are given the expressions $x^2-25=(x+5)(x-5)$,$x$, and $x-5$. Thus their LCM is $x(x+5)(x-5)$.

\section*{14}
We are given the expressions $9x^2-16=(3x+4)(3x-4)$ and $3x^2-2x-8=(3x+4)(x-2)$. Thus their LCM is $(3x+4)(3x-4)(x-2)$. 

\section*{15}
We are given the expressions $3x^2-9x=3x(x-3)$ and $6x^2=3x(2x)$. Thus their LCM is $3x(x-3)2x=6x(x-3)$. B

\section*{18}
\[ \dfrac{x-4}{5x}-\dfrac{12}{5(x-4)}=\dfrac{(x-4)(x-4)}{5x(x-4)}-\dfrac{12x}{5x(x-4)} \]
\[ =\dfrac{x^2-8x+16-12x}{5x(x-4)}=\dfrac{x^2-20x+16}{5x(x-4)} \]

\section*{21}
\[ \dfrac{9}{x-3}+\dfrac{2x}{x+1}=\dfrac{9(x+1)}{(x-3)(x+1)}+\dfrac{2x(x-3)}{(x+1)(x-3)} \]
\[ =\dfrac{9(x+1)+2x(x-3)}{(x+1)(x-3)}=\dfrac{9x+9+2x^2-6x}{(x+1)(x-3)} \]
\[ =\dfrac{2x^2+3x+9}{(x+1)(x-3)} \]

\section*{22}
\[ \dfrac{x+4}{x^2-4}-\dfrac{15}{x-2}=\dfrac{x+4}{(x+2)(x-2)}-\dfrac{15}{x-2} \]
\[ =\dfrac{x+4}{(x+2)(x-2)}-\dfrac{15(x+2)}{(x-2)(x+2)}=\dfrac{x+4-15x-30}{(x+2)(x-2)} \]
\[ =\dfrac{-14x-26}{(x+2)(x-2)} \]

\section*{23}
\[ \dfrac{-15x}{x^2-8x+16}+\dfrac{12}{x-4}=\dfrac{-15x}{(x-4)^2}+\dfrac{12}{x-4} \]
\[ =\dfrac{-15x}{(x-4)^2}+\dfrac{12(x-4)}{(x-4)^2}=\dfrac{-15x+12(x-4)}{(x-4)^2} \]
\[ =\dfrac{-3x-48}{(x-4)^2}=\dfrac{-3(x+48)}{(x-4)^2} \]

\section*{24}
\[ \dfrac{x^2-5}{x^2+5x-14}-\dfrac{x+3}{x+7}=\dfrac{x^2-5}{(x+7)(x-2)}-\dfrac{x+3}{x+7} \]
\[ =\dfrac{x^2-5}{(x+7)(x-2)}-\dfrac{(x+3)(x-2)}{(x+7)(x-2)} \]
\[ =\dfrac{x^2-5-(x+3)(x-2)}{(x+7)(x-2)}=\dfrac{x^2-5-x^2-x+6}{(x+7)(x-2)} \]
\[ =\dfrac{-x+1}{(x+7)(x-2)} \]

\section*{26}
\[ \dfrac{2x}{x+4}-\dfrac{x^2+4}{x^2-16}=\dfrac{2x}{x+4}-\dfrac{x^2+4}{(x-4)(x+4)} \]
\[ =\dfrac{2x(x-4)}{(x+4)(x-4)}-\dfrac{x^2+4}{(x-4)(x+4)} \]
\[ =\dfrac{2x^2-8x-x^2-4}{(x+4)(x-4)}=\dfrac{x^2-8x-4}{(x-4)(x+4)} \]
C

\section*{28}
\[ \dfrac{x+3}{x^2-2x-8}-\dfrac{x-5}{x^2-12x+32}=\dfrac{x+3}{(x-4)(x+2)}-\dfrac{x-5}{(x-4)(x-8)} \]
\[ =\dfrac{(x+3)(x-8)}{(x-4)(x+2)(x-8)}-\dfrac{(x-5)(x+2)}{(x-4)(x+2)(x-8)} \]
\[ =\dfrac{x^2-5x-24-(x^2-3x-10)}{(x-4)(x+2)(x-8)} \]
\[ =\dfrac{-2x-14}{(x-4)(x+2)(x-8)}=\dfrac{-2(x+7)}{(x-4)(x+2)(x-8)} \]

\section*{30}
\[ \dfrac{x+3}{x^2-25}-\dfrac{x-1}{x-5}+\dfrac{3}{x+3} \]
\[ =\dfrac{x+3}{(x+5)(x-5)}-\dfrac{x-1}{x-5}+\dfrac{3}{x+3} \]
\[ =\dfrac{(x+3)^2}{(x+5)(x-5)(x+3)}-\dfrac{(x-1)(x+3)(x+5)}{(x-5)(x+5)(x+3)}+\dfrac{3(x+5)(x-5)}{(x+3)(x-5)(x+5)} \]
\[ =\dfrac{(x^2+6x+9)-(x^2+2x-3)(x+5)+(3x^2-75)}{(x+5)(x+3)(x-5)} \]
\[ =\dfrac{4x^2+6x-66-(x^3+2x^2-3x+5x^2+10x-15)}{(x+5)(x+3)(x-5)} \]
\[ =\dfrac{-x^3-3x^2-x-51}{(x+5)(x+3)(x-5)} \]

\section*{31}
\[ \dfrac{\dfrac{x}{3}-6}{10+\dfrac{4}{x}}=\dfrac{\dfrac{x-18}{3}}{\dfrac{10x+4}{x}} \]
\[ =\dfrac{x^2-18x}{30x+12} \]

\section*{32}
\[ \dfrac{15-\dfrac{2}{x}}{\dfrac{x}{5}+4}=\dfrac{\dfrac{15x-2}{x}}{\dfrac{x+20}{5}} \]
\[ =\dfrac{75x-10}{x^2+20x} \]

\section*{33}
\[ \dfrac{\dfrac{16}{x-2}}{\dfrac{4}{x+1}+\dfrac{6}{x}} \]
\[ =\dfrac{\dfrac{16}{x-2}}{\dfrac{10x+6}{x(x+1)}} \]
\[ =\dfrac{16x(x+1)}{10x^2-14x-12}=\dfrac{16x(x+1)}{2(5x^2-7x-6)}=\dfrac{8x(x+1)}{5x^2-7x-6} \]
\[ =\dfrac{8x(x+1)}{(5x+3)(x-2)} \]

\section*{35}
\[ \dfrac{\dfrac{3}{x-2}-\dfrac{6}{x^2-4}}{\dfrac{3}{x+2}+\dfrac{1}{x-2}} \]
\[ =\dfrac{\dfrac{3}{x-2}-\dfrac{6}{(x+2)(x-2)}}{\dfrac{3}{x+2}+\dfrac{1}{x-2}} \]
\[ =\dfrac{\dfrac{3}{x-2}-\dfrac{6}{(x+2)(x-2)}}{\dfrac{3}{x+2}+\dfrac{1}{x-2}}\times\dfrac{(x+2)(x-2)}{(x+2)(x-2)} \]
\[ =\dfrac{3(x+2)-6}{3(x-2)+(x+2)} \]
\[ =\dfrac{3x}{4x-4} \]

\section*{36}
\[ \dfrac{\dfrac{1}{3x^2-3}}{\dfrac{5}{x+1}-\dfrac{x+4}{x^2-3x-4}} \]
\[ =\dfrac{\dfrac{1}{3(x+1)(x-1)}}{\dfrac{5}{x+1}-\dfrac{x+4}{(x-4)(x+1)}} \]
\[ =\dfrac{\dfrac{1}{3(x+1)(x-1)}}{\dfrac{5}{x+1}-\dfrac{x+4}{(x-4)(x+1)}}\times\dfrac{3(x-4)(x+1)(x-1)}{3(x-4)(x+1)(x-1)} \]
\[ =\dfrac{x-4}{15(x-4)(x-1)-3(x+4)(x-1)} \]
\[ =\dfrac{x-4}{15x^2-75x+60-(3x^2+9x-12)} \]
\[ =\dfrac{x-4}{12x^2-84x+72} \]
\[ =\dfrac{x-4}{12(x^2-7x+6)} \]
\[ =\dfrac{x-4}{12(x-6)(x-1)} \]

\end{document}
